% =========================================================================
%  CHAPTER 18: GALACTIC STRUCTURE
%  Volume III: The Consequences — Part G
% =========================================================================
%  STATUS: COMPLETE
%  SOURCE: Extracted from Book_3/Chapter_10_Cosmology.tex §2
% =========================================================================

\chapter{Galactic Structure: The Dynamic Pressure Vortex}
\label{ch:galactic-structure}


% =============================================
\section{The Problem of Flat Rotation Curves}
\label{sec:rotation-curves}
% =============================================

The most consequential unsolved problem in gravitational dynamics is the flat rotation curve of spiral galaxies.  Newtonian gravity predicts that stellar orbital velocities should decline as $v \propto r^{-1/2}$ beyond the luminous disc --- the same Keplerian fall-off observed perfectly within the solar system.  Instead, observed velocities remain approximately constant far beyond the visible edge of the galaxy.

The standard resolution invokes \textbf{dark matter}: an invisible, non-baryonic substance that forms a massive halo around each galaxy, providing the additional gravitational pull required to explain the rotation curve.  Despite decades of direct searches, dark matter has never been detected.

SDT offers a purely mechanical alternative.


% =============================================
\section{The Mechanism: Retarded Pressure Fields}
\label{sec:retarded-field}
% =============================================

A galaxy is not a static collection of stars held by an invisible scaffolding.  It is a \textbf{rotating displacement vortex} governed by the dynamic, relativistic effects of its own propagating pressure field.

\subsection{The Spiral Pressure Wave}

The galaxy's rotating central mass creates a \textbf{spiral pressure wave} propagating outwards at $c$.  This wave is not an instantaneous field; it is a physical pressure disturbance that takes finite time to reach outer stars.

\begin{enumerate}
  \item The rotating nucleus generates a time-varying pressure field.
  \item This field propagates as a spiral wave at $c$.
  \item Outer stars receive a continuous, non-zero \textbf{tangential force} from the retarded field.
  \item This tangential component transfers angular momentum from inner to outer regions.
\end{enumerate}

\subsection{The SDT Velocity Profile}

\begin{equation}\label{eq:v-galaxy}
  v_{\text{SDT}}(r) = \sqrt{r \cdot a_r(r)}
\end{equation}

where $a_r$ is the radial acceleration from the full, retarded potential of the galaxy's baryonic mass.  The critical difference from Newtonian gravity: the retarded field includes contributions that the instantaneous approximation misses.

\textbf{Consequence:} The continuous energy input from the spiral pressure wave prevents Keplerian velocity decay.  The ``missing mass'' of dark matter is the missing \emph{physics} of a static gravitational model applied to a dynamic, relativistic system.


% =============================================
\section{Comparison with Observations}
\label{sec:rotation-comparison}
% =============================================

The retarded pressure profile naturally produces:
\begin{itemize}
  \item \textbf{Flat asymptotic velocity:} The tangential component of the spiral wave provides a velocity floor.
  \item \textbf{Rising curve in the inner disc:} Near the nucleus, the field is nearly instantaneous and follows Keplerian dynamics.
  \item \textbf{Correlation with baryonic mass:} The Tully--Fisher relation ($v^4 \propto L$) emerges naturally from the $\kop$-based scaling: galaxies with more baryonic mass have larger $c/\kop_{\text{gal}}$ values.
\end{itemize}

\textbf{No dark matter is required.}  The ``missing mass'' is the retarded contribution to the pressure field --- real, physical, and entirely baryonic.


% =============================================
\section{Summary}
% =============================================

\begin{enumerate}
  \item Galaxies are rotating displacement vortices, not static mass distributions.
  \item The retarded (finite-propagation-speed) pressure field includes a tangential component that transfers angular momentum outward.
  \item This mechanical effect produces flat rotation curves without invoking dark matter.
  \item The Tully--Fisher relation emerges naturally from the $\kop$-based scaling.
\end{enumerate}
